\section{A paper proof of the auxiliary-prime selection}\label{pr:sec}

Fix \(0<\alpha\le1/4,\ K\ge1\). Write \(L=\log Y,\ Z=Y^{400K}\).
Consider the primes \(\mathcal P\) between \(\sqrt U\) (excluded) and \(U\)
in a reduced class \(a\bmod m\), with \(8\mid m\). Assume
\[
 L^2\log m\le4\log U,\qquad
 Y^{\alpha/6}\le\log U\le Y^{\alpha/5}L^2 .
\]
Every prime divisor of \(m\) is at most \(Y^{\alpha/5}\) or greater than \(Z\).
All thresholds below are uniform in \(m,a,U\) subject to these inequalities.
We prove the estimates on a whole interval used in Section~\ref{sup:section}.

\begin{lemma}[Primes and quadratic signs in a progression]\label{pr:signs}
For large \(Y\), \(|\mathcal P|\ge U^{9/10}\). For each fixed \(A,B>0\)
there is at most one exceptional squarefree kernel \(s_*\) to be omitted
from the following assertion. For odd squarefree \(1<s\le Y^A\) coprime to \(m\),
\[
 \left|\sum_{p\in\mathcal P}(p/s)\right|\le Y^{-B}|\mathcal P| .
\]
\end{lemma}
\begin{proof}
For the count, \(m\le U^{1/32}\) eventually. Let \(D=2^j m\) be the largest
such multiple at most \(U^{1/32}\). Then \(U^{1/32}/2<D\le U^{1/32}\);
\(a\) remains coprime to \(D\). Maynard's unconditional Theorem~3.2
\cite{maynard} gives
\[
 \pi(U;D,a)\gg \frac{\log D}{\log U}\frac{U}{\sqrt D\,\varphi(D)}
              \gg U^{61/64}.
\]
The theorem requires \(D\) sufficiently large and \(U\ge D^8\); both follow
from the construction. At most \(\sqrt U\) primes have been counted below the
desired interval. The fixed positive constant is absorbed by the strict
inequality \(61/64>9/10\).

For the second assertion put \(D_*=mY^A\). Uniformly for \(t\in[\sqrt U,U]\)
and \(M=ms\), one has \(t\ge M^{12}\) eventually and
\(\log t/\log D_*\gg L^2\). Indeed \(\log D_*=\log m+AL\),
and \(Y^{\alpha/6}/L^3\to\infty\). Apply Thorner--Zaman, Corollary~1.4
and Remark~1.5 \cite{tz}, to each unit class \(b\bmod M\). In terms of
\(\Theta_b(t)=\sum_{p\le t,\ p\equiv b\bmod M}\log p\), this gives
\[
 \Theta_b(t)=H_b(t)(1+O(\Delta)),\qquad
 H_b(t)=\frac{t-\chi(b)t^\beta/\beta}{\varphi(M)} .
\]
The second term is omitted when the exceptional real zero of modulus \(M\)
is absent. The main terms are positive in this range, by (1.7) of
\cite{tz}. The relative error is
\(\Delta=\exp(-c\log t/\log M)+
 \exp(-c(\log t)^{3/5}/(\log\log t)^{1/5})\).
For any fixed \(B'>0\), this is \(O(Y^{-B'})\) uniformly here.

There is an absolute \(c_0>0\) such that at most one primitive character
of conductor at most \(D_*\) can have a real zero in
\((1-c_0/\log D_*,1)\). Indeed two would induce distinct characters
modulo the least common multiple of their conductors, at most \(D_*^2\).
The classical uniqueness of the exceptional zero for that one modulus
(the zero-free region recalled in \cite[\S1.1]{tz}) excludes this when
\(c_0\) is sufficiently small. The extra Euler factors of an induced
character do not vanish on this interval. If this character exists, take
\(s_*\) to be the product of the primes of its conductor not dividing \(m\).

Sum the displayed formula over classes \(b\equiv a\bmod m\) in
\((\mathbf Z/M)^\times\), first without weights, then with weights
\(\lambda(b)=(b/s)\). The Chinese remainder theorem makes the second
sum of the primary terms zero. The exceptional term in the weighted sum
also vanishes unless the component of \(\chi\) modulo \(s\) equals
\(\lambda\): the product of two distinct real characters has sum zero.
In this last case that component is nontrivial at every prime of \(s\).
Its primitive conductor outside \(m\) has radical \(s\). For \(s\ne s_*\)
we must then have \(1-\beta\ge c_0/\log D_*\). The \emph{unweighted}
exceptional term sums to zero in this case, and the magnitude of the
weighted main term divided by the unweighted main sum is
\(t^{\beta-1}/\beta=O(\exp(-c_0\log t/\log D_*))\);
the exceptional zeros under consideration have \(\beta\ge1/2\).
In every case the absolute error summed over classes is
\(O(\Delta\sum_b H_b(t))\), since the individual estimates are relative
to positive main terms. This remains true when the base progression
itself has a small main term.

It follows that the signed prefix sum of \(\log p\), on these unit classes,
is \(O(Y^{-B'})\) times the unsigned prefix sum for \(a\bmod m\).
The latter is at most \((|\mathcal P|+\sqrt U)\log U\) throughout
\([\sqrt U,U]\). Partial summation with \(1/\log t\) on this interval
costs \(O(1/\log U)\), including both endpoints. Primes dividing \(s\)
are below \(\sqrt U\). Use the first assertion to absorb \(\sqrt U\)
and take \(B'>B\).
\end{proof}

\begin{proposition}[A uniform mean square]\label{pr:ms}
For \(|\kappa(n)|\le1\), put
\[
 S_p(y)=\sum_{n\le Z}\Lambda(n)\kappa(n)(n/p)\min(1,y/n).
\]
For \(Y\le y\le Y^K\) one has
\(\sum_{p\in\mathcal P}|S_p(y)|^2\ll_K |\mathcal P|y\log y\).
For each fixed \(\tau>0\), the proportion of primes for which
\(|S_p(y)|>\tfrac12y^{1/2+\tau}\) somewhere in this interval is
\(O_{K,\tau}(Y^{-\tau})\), uniformly in \(\kappa\).
\end{proposition}
\begin{proof}
Eventually \(\sqrt U>Z^2\). For \(n=\ell^j\), let \(r(n)=\ell\) if
\(j\) is odd and \(\ell\nmid m\), and \(r(n)=1\) otherwise. Quadratic
reciprocity expresses \((n/p)\) as a fixed sign times \((p/r(n))\):
the signs at \(2\) and at primes of \(m\) are fixed by the progression,
and \(r(n)\) is odd. Group the absolute coefficients by their square
class, writing \(b_g=\sum_{r(n)=g,\ n\le Z}\Lambda(n)\min(1,y/n)\)
in the finite elementary two-group generated by the primes up to \(Z\)
outside \(m\).

The elementary Chebyshev bound \(\sum_{n\le u}\Lambda(n)\ll u\)
gives \(\sum_g b_g\ll y(1+\log Z)\) by partial summation.
Even powers have cumulative weight at most \(C\sqrt u\) up to \(u\);
integration against \(y\,du/u^2\) for \(u\ge y\) bounds their ramp
weight by \(2C\sqrt y\). For each prime \(\ell\le y\), all its powers
together contribute at most \(\log y+2\log\ell\): after the last power
at most \(y\) the geometric series has total at most \(2\).
Thus \(b_1\ll\sqrt y+Y^{\alpha/5}\log y\ll\sqrt y\).
Among other groups with \(\ell\le y\), at most \(O(y/\log y)\) have
positive weight, each of size \(O(\log y)\). For \(\ell>y\) the bound
is \(2y\log\ell/\ell\). On each dyadic block above \(y\) Chebyshev's
bound gives \(\sum_{\ell>y}(\log\ell)^2/\ell^2\ll(\log y)/y\).
Consequently \(\sum_g b_g^2\ll y\log y\).

In the expansion of the second moment the square class \(gh\) is odd,
squarefree, coprime to \(m\) and at most \(Z^2\) whenever the coefficient
is nonzero. Its normalized prime sum costs at most \(Y^{-K-1}\), except
for \(gh=1\) and possibly \(gh=s_*\), by Lemma~\ref{pr:signs}.
For either exceptional value \(d\), multiplication by \(d\) permutes
the classes; Cauchy--Schwarz gives \(\sum_g b_g b_{gd}\le\sum_g b_g^2\).
The remaining total is \(O(Y^{-K-1}y^2(1+\log Z)^2)\ll_K y\log y\).
This proves the pointwise estimate.

Except at integer breakpoints, \(y S'_p(y)\) is the same sum with
\(\kappa(n)\) replaced by \(\kappa(n)\mathbf1_{n>y}\). Therefore
\(\sum_p|S'_p(y)|^2\ll_K|\mathcal P|(\log y)/y\).
On a dyadic block \([u,v]\subseteq[Y,Y^K]\), \(v\le2u\), integration
of the derivative and Cauchy--Schwarz bound
\(\sum_p\sup_{[u,v]}|S_p|^2\ll_K|\mathcal P|u\log(2u)\).
The proportion violating the asserted threshold on the block is
\(O_K(u^{-2\tau}\log(2u))=O_{K,\tau}(u^{-\tau})\). Sum the geometric
series over blocks, truncating the last block at \(Y^K\). For \(K=1\)
the pointwise estimate suffices.
\end{proof}

\begin{lemma}[Truncation using a finite-height box]\label{pr:perron}
Let \(\chi\) be a primitive nonprincipal character of conductor at most
\(Q\), with \(L(s,\chi)\ne0\) for \(\Re s>399/400,\ |\Im s|\le Q\).
For \(Q\) and \(z\) sufficiently large, \(z\) a half-integer and
\(2\le\log z\le Q^{1/8}\), one has
\[
 \left|\frac{L'}{L}(1,\chi)+\sum_{n\le z}\frac{\Lambda(n)\chi(n)}n\right|
 \ll z^{-1/400}(\log z)^2(\log Q)^2+
        \frac{(\log z)^2(1+\log Q)}{Q}.
\]
\end{lemma}
\begin{proof}
Put \(\sigma_0=399/400,\ \eta=1/\log z,\ T=Q/2\). Assume
\(\eta\le1/800\), as is sufficient for the assertion. For \(|t|\le T\)
the open disk of radius \(2-\sigma_0\) about \(2+it\) is zero-free.
There is a branch of \(\log L\) there with bounded central value, by
the absolutely convergent Euler product at real part two. Periodicity
and partial summation give \(|L(s,\chi)|\le Q|s|/\Re s\) in the disk,
so \(\Re\log L\ll\log Q\). For clarity, if \(F\) is the branch minus
its central value, take \(H\gg\log Q\) with \(\Re F\le H\).
The map \(F/(2H-F)\) takes the disk into the unit disk and vanishes at
the centre. Schwarz's lemma gives \(|F|\ll H/\eta\) on the disk
shortened by \(\eta/2\); Cauchy's derivative estimate on a circle of
radius \(\eta/2\) then gives
\(|L'/L(\sigma+it,\chi)|\ll(\log Q)/\eta^2\) for
\(\sigma_0+\eta\le\sigma\le1+\eta\).

Integrate \(-\frac{L'}L(1+w,\chi)z^w/w\) up the segment
\(\Re w=\eta,\ |\Im w|\le T\), with factor \(1/(2\pi i)\).
Its difference from the prime-power sum is \(O((\log z)^2/T)\).
Indeed the scalar Perron kernel for \(v=z/n\ne1\) has error
\(O(v^\eta/(T|\log v|))\), by integration by parts on the tails of
the vertical inversion integral for \(\mathbf1_{v>1}\).
Outside \(z/2<n<2z\) use \(\sum(\log n)n^{-1-\eta}\ll\eta^{-2}\).
Inside use \(|\log(z/n)|\ge|z-n|/(2z)\), \(|z-n|\ge1/2\), and sum
the resulting harmonic series. Absolute convergence on the initial
line justifies interchanging the sum and the finite integral.
Shift to \(\Re w=-1/400+\eta\); the only residue is
\(-L'(1,\chi)/L(1,\chi)\) at zero. The two horizontal integrals cost
\(O((\log Q)\eta^{-2}/T)\). On the new vertical line
\(|z^w|\le e z^{-1/400}\) and \(\int_{-T}^T|w|^{-1}dt\ll\log(2T)\),
since \(1/400-\eta\ge1/800\). The derivative bound finishes the proof.
\end{proof}

\begin{theorem}[Uniform selection of the two primes]\label{pr:pair}
Fix \(M\ge1,\ 0<\varepsilon<1/52,\ K\ge1\). For every sufficiently
large \(Y\), put \(z_0=Y^{\varepsilon/5}\), \(z_1=z_0(\log Y)^2\).
There are primes \(q\in(e^{z_0/2},e^{z_0}],\
d\in(e^{z_1/2},e^{z_1}]\) with \(q\equiv3\bmod8,\ d\equiv7\bmod16\)
and \(d\equiv-1\) modulo \(q\) and every odd prime at most \(z_0\).
They are coprime to \(8M\). For every primitive inducer \(\xi\) of a
character modulo \(8M\), and every \((i,j)\in\{0,1\}^2\setminus\{(0,0)\}\),
write \(\chi\) for the primitive character associated to
\(\xi(\cdot)(\cdot/q)^i(\cdot/d)^j\). Simultaneously for these choices,
\[
 |G_\chi(y)|\le y^{1/2+\varepsilon/2},\qquad
 |L'(1,\chi)/L(1,\chi)|\le4\log y\qquad(Y\le y\le Y^K).
\]
Every selected twist is nonprincipal and has the zero-free box in
Lemma~\ref{pr:perron} up to the conductor bound used at its selection.
\end{theorem}
\begin{proof}
For one progression take a list of at most \(R\) primitive base characters
of conductor at most \(F\le\sqrt U\), where \(R\) is fixed. For \(p\in
\mathcal P\) twisting a base character of conductor \(f\) by \((\cdot/p)\)
gives primitive conductor \(fp\), by the prime-by-prime conductor factorization.
For \(Q=FU\) the density estimate of Lemma~\ref{sup:density} bounds the
number of bad primes by \(O(Q^{1/4})\): a chosen bad character determines
its prime uniquely because its conductor has only one prime factor \(>\sqrt U\).
Divide by \(|\mathcal P|\ge U^{9/10}\) and use \(Q\le U^{3/2}\).
By Proposition~\ref{pr:ms} a further \(O_{K,\tau}(RY^{-\tau})\) proportion
fails the ramp estimate for some base character. Both proportions tend to zero.

For any survivor choose \(z=\lfloor Y^{400K}\rfloor+1/2\) in
Lemma~\ref{pr:perron}. Since \(\log Q\ll Y^{\alpha/5}(\log Y)^2\),
the error is \(O_K(Y^{-K+2\alpha/5}(\log Y)^6)\). The term divided by \(Q\)
is smaller, since \(\log Q\ge Y^{\alpha/6}\). The elementary identity
\[
 G_\chi(y)=-S_p(y)+y\left(\frac{L'}L(1,\chi)
                           +\sum_{n\le z}\frac{\Lambda(n)\chi(n)}n\right)
\]
follows from \(y\ge1,\ y\le z\) by splitting the ramp at \(n=y\).
The half-integer retains exactly the terms \(n\le Y^{400K}\), and
\(y\le Y^K<z\) for large \(Y\).
For \(\tau=\varepsilon/2\) and either \(\alpha=\varepsilon\) or
\(6\varepsilon/5\), the tail after multiplication by \(y\) is
\(o(Y^{1/2+\tau})\), giving the asserted ramp bound.
For the logarithmic derivative take instead
\(z=\lfloor(\log Q)^{1000}\rfloor+1/2\).
The error in Lemma~\ref{pr:perron} is \(o(1)\). The elementary estimate
\(\sum_{n\le u}\Lambda(n)/n=\log u+O(1)\) follows from
\(\sum_{n\le u}\Lambda(n)\lfloor u/n\rfloor=\log(\lfloor u\rfloor!)\),
Chebyshev's bound and Stirling's formula. Hence the derivative is bounded
by \(1000\log\log Q+O(1)\).

First use \(m=8,\ a=3,\ U=e^{z_0},\ \alpha=\varepsilon,\ F=8M\)
and all primitive inducers from \(8M\). Fix a good \(q\). Next take
\[
 m=16q\prod_{\substack{\ell\le z_0\\\ell\ {\rm odd\ prime}}}\ell,\quad
 U=e^{z_1},\quad \alpha=6\varepsilon/5,\quad F=(8M)e^{z_0}.
\]
The Chinese remainder theorem gives the required reduced class \(a\).
For large \(Y\), \(q>Z,\ F\le\sqrt U,\ Y^{\alpha/6}\le z_1
\le Y^{\alpha/5}L^2\). Also \(\log m\le4z_0\): the elementary bound
\(\sum_{\ell\le u}\log\ell\le(2\log2+o(1))u\) follows by bounding
the primes between \(n\) and \(2n\) by the middle binomial coefficient
and summing dyadically. Thus \(L^2\log m\le4z_1\).
Use both the old characters and their \(q\)-twists, at most \(2\varphi(8M)\)
characters. The resulting \(d\) provides the two remaining twist patterns.
At both stages \(\log Q\le3Y^{\varepsilon/5}L^2\) eventually, so
\(1000\log\log Q\le200\varepsilon L+O(\log L)\). The strict inequality
\(200\varepsilon<4\) gives the promised bound. All choices of \(q,d\)
precede the choice of \(y\).
\end{proof}
